\subsection{Shapley-Shubik Power Index}

\content{
        \begin{definition} A sequential coalition is an ordered coalition, that
        is, it keeps track of the order in which each player joins the
        coalition. 
        \end{definition}
}{% presentation
}{% nopresentation
}{% comments
We will use the notation 
$$ \la P_1, P_2, P_3\ra$$
to denote sequential coalitions. 
}

\content{
        \begin{definition} A \textbf{pivotal player} is the player in the
        coalition which changes the coalition from a losing coalition to a
        winning one.
        \end{definition}
}{% presentation
}{% nopresentation
}{% comments
}

\content{
        \begin{definition} (Shapley-Shubik power index) To compute the
        Shapley-Shubik power index: 
        \begin{itemize}
        \item List all sequential coalitions
        \item Determine all pivotal players
        \item Count up how many times each player was pivotal
        \item convert these into fractions by dividing each count by the total
        number of sequential coalitions
        \end{itemize}
        \end{definition}
}{% presentation
}{% nopresentation
}{% comments
}

\content{
        \begin{example} Calculate the Shapley-Shubik Power Index for the following
        weighted voting system
        $$ [ 8: 6,3,2 ]. $$
        \end{example}
}{% presentation
\vspace{3in}
}{% nopresentation
\vspace{3in}
}{% comments
Power index is: 
$$ P_1: 2/3, P_2: 1/6, P_3: 1/6.$$
}


\content{
        \begin{example} Compute the Shapley-Shubik Power Index for the weighted
        voting system: 
        $$ [ 36 : 20, 17, 15 ] $$
        \end{example}
}{% presentation
\vspace{3in}
}{% nopresentation
\vspace{3in}
}{% comments
$P_1 = 1/2$, $P_2 = 1/2$, $P_3 = 0$.
}

\content{
        \begin{example} Compute the Shapley-Shubik Power Index for the weighted
        voting system: 
        $$ [ 22: 23, 15, 7 ]. $$
        \end{example}
}{% presentation
\vspace{2in}
}{% nopresentation
\vspace{2in}
}{% comments
Notice that $P_1$ is \textbf{not} a dictator!
}

\content{
        \begin{example} Small review: Compute the Banzhaf Power Index for the weighted
        voting system: 
        $$ [ 39: 23, 17, 15 ]. $$
        Are there any dummies in this system? 
        \end{example}
}{% presentation
\vspace{3in}
}{% nopresentation
\vspace{3in}
}{% comments
Power index: 
$$ P_1: 1/2, P_2: 1/2, P_3: 0. $$
}

\content{
        \begin{example} Consider the weighted voting system given by 
        $$ [ q: 7, 3, 1 ]. $$
        Which values of $q$ result in a dictator? (list all possible values)
        \end{example}
}{% presentation
\vspace{3in}
}{% nopresentation
\vspace{1in}
}{% comments
}

\content{
        \begin{example} In the Shapley-Shubik method, is it possible for a dummy
        to be a pivotal player? Why or why not? 
        \end{example}
}{% presentation
\vspace{2in}
}{% nopresentation
\vspace{2in}
}{% comments
No, if they were, they would be a critical player in a winning coalition. 
}
